\section{A clock-free law on late states}\label{sec:clockfree}
Age-dependent cone separation proves Theorem~\ref{thm:realization},
but does not produce a fixed statistic of a late state. We localize
error certificates with a seed-independent fault envelope.
Independence belongs to local certificate events, not actual errors.

\subsection{Independent local certificates}
Implement updates with independent uniforms $U_{z,t}$, putting the
deterministic output first in the sampling interval. Actual faults
are contained in the independent envelope
\[
 {\cal F}=\{(z,t):U_{z,t}\geq1-\varepsilon_*\}.
\]
It is independent of the seed. Put
\[
 D_*=24,\quad c=16D_*^8\varepsilon_*=\tfrac14,\quad
 \delta_*=(1+8D_*)\varepsilon_*\leq\tfrac1{16}.
\]
G\'acs's explanation-tree and weighted-tree counting lemmas
\cite[Lemmas~2.3 and 2.6]{gacstoom}, in the PDF numbering, give:
a nonfaulty error has an explanation with $n$ fault vertices and
at most $4(n-1)$ edges; the graph degree is $D_*$; each edge moves
each space-time coordinate by at most one; and the number of
weighted rooted trees of $\ell$ edges is at most
$2D_*(2D_*^2)^\ell$.

\begin{lemma}[Certificate domination]\label{lem:certificates}
At $m$ equal-time endpoints of pairwise sup-distance at least
$L\geq8$, for $1\leq h\leq m$ and every cloned row,
\begin{equation}\label{eq:certificate}
 \Pr(\text{at least }h\text{ errors}\mid Y)
 \leq8D_*m\varepsilon_*c^{\lceil L/8\rceil-1}
                                  +\binom mh\delta_*^h.
\end{equation}
The estimate is uniform in time and $K\in U_\beta$.
\end{lemma}
\begin{proof}
With $A_*=2D_*^2$, the number of allowed weighted trees with at
most $4(n-1)$ edges is bounded by
$2D_*\sum_{\ell=0}^{4(n-1)}A_*^\ell
\leq4D_*A_*^{4(n-1)}$. Requiring their $n$ marked vertices
to lie in the fault envelope has probability $\varepsilon_*^n$;
extra faults are allowed. The union $S_n(u)$ of such certificates
at a root therefore satisfies
$\Pr(S_n(u))\leq4D_*\varepsilon_*c^{n-1}$.
Each tree lies in a sup-ball of radius $4(n-1)$.

Let $n_0=\lceil L/8\rceil$. At endpoint $j$ let $A_j$ be the
root-fault event or a certificate with $n<n_0$. Its supporting
coins lie at radius strictly less than $L/2$. These supports are
disjoint for distinct endpoints. The $A_j$ are independent,
including conditional on the row, and
\[
 \Pr(A_j)\leq\varepsilon_*+
       4D_*\varepsilon_* /(1-c)\leq\delta_*.
\]
The remaining large-certificate event has probability at most
$4D_*\varepsilon_*c^{n_0-1}/(1-c)
\leq8D_*\varepsilon_*c^{n_0-1}$ per root.
Every actual error belongs to a small or large event. Union-bound
the large events and the choices of $h$ independent small events.
This gives \eqref{eq:certificate}. Intersecting supports with the
available time half-space only reduces them.
\end{proof}

Actual errors can remain correlated at every transverse distance.
For a fair cloned bit $Y$, take time-one error probabilities
$p_0,p_1$ depending on its value. Distant errors are conditionally
independent but have unconditional covariance $(p_1-p_0)^2/4$.
The lemma does not claim decay of this covariance.

\subsection{Scale-only geometry and cluster states}
At scale $k\geq1$ choose
\begin{equation}\label{eq:scalegeometry}
 m_k=8k+9,\qquad L_k=64(k+1),\qquad h_k=2^{-k}.
\end{equation}
Decode each blocked position by strict majority of $m_k$ copies
spaced by $L_k$, then unpack. Write $r_k$ for the resulting gap
ratio. The two words use at most $2k+3$ blocks. With $e_k$ the
right side of \eqref{eq:certificate} at majority threshold,
\begin{align}
 e_k&\leq(8k+9)2^{-16k-14}+2^{-8k-10}
                       \leq(8k+10)2^{-8k-10},\label{eq:ek}\\
 (2k+3)e_k&\leq90\,2^{-6k-12}\leq4^{-k}/96.\label{eq:budget}
\end{align}
Indeed $8D_*\varepsilon_*\leq1$, $c=1/4$ and $\delta_*\leq1/16$.
Also $(2k+3)(8k+10)\leq90\,4^{k-1}$: equality holds at $k=1$
and successive polynomial ratios are less than four. The last
comparison holds at $k=1$ and improves by a factor sixteen.
These are all-scale bounds, not numerical fits.

Conditional on the row, apply the lemma at majority threshold,
union-bound over the word's blocks, and average over the row.
For either gap word $u$ and every time,
\begin{equation}\label{eq:fixedworderror}
 \left|\Pr_{\lambda P_K^t}(\text{scale-}k\text{ decoded }u)
             -(F_*^t\mu_{\rm row})([u])\right|\leq\eta_k:=4^{-k}/96.
\end{equation}
The reference probabilities converge. Finite decoded probabilities
are continuous affine functions of the physical state, so every
ordinary cluster state inherits their $\eta_k$ bounds against
$\nu_b$. Ces\`aro cluster states inherit them too, by averaging
the convergent reference sequence.

For nonnegative pairs with $s=a+b>0$ and individual perturbations
at most $\eta\leq s/4$, subtraction gives
\begin{equation}\label{eq:ratiobound}
 \left|\frac{a'}{a'+b'}-\frac a{a+b}\right|
 \leq\frac{\eta}{s-2\eta}\leq\frac{2\eta}{s}.
\end{equation}
The numerator is $b(a'-a)-a(b'-b)$. Since $s\geq2^{-k}/6$,
every ordinary or Ces\`aro cluster state satisfies
\begin{equation}\label{eq:cluster-ratio}
 |r_k(\bar\mu)-b_k|\leq h_k/8.
\end{equation}
The decoded denominator is positive and $\lim_k r_k(\bar\mu)=\beta$.
This alone does not give fixed-$k$ inner time limits.

\subsection{Continuous quantization gives genuine inner limits}
Choose the realization with
$b_k=2^{-k}\lfloor2^ka_k\rfloor$, for uniformly $0'$-computable
rationals $a_k\in[0,1]$ tending to $\beta$. This is exact rational
flooring and $b_k\to\beta$. On $[0,1]$ define
\begin{equation}\label{eq:quantizer}
 Q_k(x)=h_k\sum_{j=0}^{2^k-1}
   \min\{1,\max\{0,2x/h_k-2j-\tfrac12\}\}.
\end{equation}
This target-independent continuous rational piecewise-linear map
equals $jh_k$ on $|x-jh_k|\leq h_k/4$, intersected with $[0,1]$,
including the endpoints. Use $s_k=Q_k\circ r_k$.

\begin{proof}[Proof of Theorem~\ref{thm:clockfree}]
Equation~\eqref{eq:cluster-ratio} gives $s_k(\bar\mu)=b_k$
exactly at every cluster state. At fixed $k$ the same estimate gives
$\limsup_t|r_k(\lambda P_K^t)-b_k|\leq h_k/8$.
There is a strict margin to the plateau boundary $h_k/4$.
Hence $s_k(\lambda P_K^t)=b_k$ for all sufficiently large $t$.
Every inner time limit exists and the outer limit is $\beta$,
uniformly on the kernel region. Denominators are positive at
finite times by positive rates, and at cluster states by the
displayed bounds. The statistic uses scale but no age. This
chooses a dyadic-parameter realization; it does not prove convergence
of every earlier unquantized ratio.

At computable $K$ and canonical seed, $C_T$ is uniformly computable
and its cluster class $V'$ has the valid-code description
\begin{equation}\label{eq:cluster-pi2}
 \mu\in V'\ \Longleftrightarrow\
 (\forall n,T)(\exists t\geq T)\ d(C_t,\mu)<2^{-n}.
\end{equation}
Strict inequality has a finite-precision positive witness, giving
a $\Pi^0_2$ condition. Compactness gives nonemptiness; TI and
\eqref{eq:cesaro-residual} give stationarity. On $V'$ the
clock-free law is $\beta$. A low name $X$ would compute that
one-limit value in $X'\equiv_T0'$, impossible if $\beta\notin\D_2$.
Every nonempty effective $\Sigma^0_2$ class has a nonempty closed
piece and a low code, so cannot have this depth-one consensus.
Ordinary cluster states share the law but need not be stationary.
\end{proof}
